% Source printed pp. 308-320 / PDF pp. 316-328.
\section{拟凝聚层与上同调理论}\label{src:7.1}
\begin{epigraph}
我真想掐死写这本书的人.
\epigraphauthor{Hermann Weyl}
\end{epigraph}
我们定义代数叠上拟凝聚层的方法, 与 \S\bookref{5.1} 中定义 Deligne--Mumford 叠上拟凝聚层的方法基本相同, 但使用 $\mathcal X$ 上的光滑平展景来代替小平展景.
\subsection{层与\texorpdfstring{$\mathcal O_{\mathcal X}$}{OX}-模}\label{src:7.1.1-subsection}
适合在一般代数叠上定义层的最小景是光滑平展拓扑. 这个景有一些技术问题, 最显著的是缺乏函子性; 见警告~\bookref{7.1.5}.
\begin{definition}[光滑平展景]\label{src:7.1.1}
代数叠 $\mathcal X$ 上的\newterm{光滑平展景}{lisse-étale site} $\mathcal X_{\text{lis-ét}}$ 是 $\mathcal X$ 上光滑概形组成的范畴, 其中态射是 $\mathcal X$ 上任意的概形态射. 覆盖 $\{U_i\to U\}$ 是一族态射, 使得 $\coprod_iU_i\to U$ 满射且平展.
\end{definition}
以 $\mathsf{Sh}(\mathcal X_{\text{lis-ét}})$ (分别地, $\mathsf{Ab}(\mathcal X_{\text{lis-ét}})$) 记 $\mathcal X_{\text{lis-ét}}$ 上层 (分别地, Abel 群层) 的范畴. 结构层 $\mathcal O_{\mathcal X}$ 定义为 $\mathcal O_{\mathcal X}(U)=\Gamma(U,\mathcal O_U)$, 从而可以定义 $\mathcal O_{\mathcal X}$-模. 对 $F\in\mathsf{Ab}(\mathcal X_{\text{lis-ét}})$, 可对任意代数叠的光滑态射 $\mathcal U\to\mathcal X$ 定义截面
\[
F(\mathcal U\to\mathcal X)=\operatorname{Eq}\bigl(F(U\to\mathcal X)\rightrightarrows F(R\to\mathcal X)\bigr),
\]
其中 $U\to\mathcal U$ 是光滑呈示, $R\to U\times_{\mathcal U}U$ 是由概形给出的平展呈示. 与 \S\bookref{5.1.1-subsection} 中 Deligne--Mumford 叠的情形一样, 这与呈示的选择无关. 特别地, 整体截面 $\Gamma(\mathcal X,F):=F(\mathcal X\xrightarrow{\mathrm{id}}\mathcal X)$ 良定义.
给定代数叠的态射 $f:\mathcal X\to\mathcal Y$, 有伴随函子
\[
\begin{tikzcd}[ampersand replacement=\&,column sep=large]
\mathsf{Sh}(\mathcal X_{\text{lis-ét}})\arrow[r,shift left,"{f_*}"]\&\mathsf{Sh}(\mathcal Y_{\text{lis-ét}})\arrow[l,shift left,"{f^{-1}}"]
\end{tikzcd}
\qquad
\begin{tikzcd}[ampersand replacement=\&,column sep=large]
\mathsf{Ab}(\mathcal X_{\text{lis-ét}})\arrow[r,shift left,"{f_*}"]\&\mathsf{Ab}(\mathcal Y_{\text{lis-ét}})\arrow[l,shift left,"{f^{-1}}"] .
\end{tikzcd}
\]
其中推前定义为 $f_*F(V\to\mathcal Y):=F(\mathcal X\times_{\mathcal Y}V\to\mathcal X)$, 逆像 $f^{-1}G$ 是预层
\begin{equation}\label{src:7.1.2}
(f^{-1}G)^{\mathrm{pre}}:(U\to\mathcal X)\longmapsto\operatorname{colim}_{(U\to V)}G(V\to\mathcal Y)
\end{equation}
的层化, 此处余极限遍历 $f$ 上的态射 $U\to V$ 的范畴, 其中 $V$ 是光滑 $\mathcal Y$-概形. 见 \cite[Tag 00WX]{SP}, 其中对景的任意连续函子建立了更一般的结论.
\begin{exercise}\label{src:7.1.3}
若 $f:\mathcal X\to\mathcal Y$ 是代数叠的光滑态射, $F$ 是 $\mathcal Y_{\text{lis-ét}}$ 上的层, 则 $f^{-1}F=F|_{\mathcal X}$, 且 $f^{-1}$ 正合.
\end{exercise}
\begin{exercise}\label{src:7.1.4}
设 $f:X\to Y$ 是概形态射.
\begin{enumerate}[(a)]
\item 设 $Y'$ 是光滑 $Y$-概形, $h_{Y'}$ 是 $Y_{\text{lis-ét}}$ 上由 $Y'$ 表示的层, 即 $h_{Y'}(V\to Y)=\operatorname{Mor}_Y(V,Y')$. 证明 $f^{-1}h_{Y'}=h_{X\times_Y Y'}$.
\item 证明 $f^{-1}\mathcal O_{Y_{\text{lis-ét}}}=\mathcal O_{X_{\text{lis-ét}}}$.

\textit{提示:} 使用 (a) 及 $\mathcal O_{Y_{\text{lis-ét}}}$ 由 $\mathbb A^1_Y$ 表示这一事实.
\item 推出若 $i_x:\Spec\kappa(x)\to X$ 是点 $x\in X$ 的嵌入, 则 $i_x^{-1}\mathcal O_{X_{\text{lis-ét}}}=\kappa(x)$. (这应与 $i_x^{-1}\mathcal O_{X_{\mathrm{Zar}}}=\mathcal O_{X,x}$ 是局部环, 以及 $i_x^{-1}\mathcal O_{X_{\text{ét}}}=\mathcal O^h_{X,x}$ 是其 Hensel 化这两个事实相比较.)
\item 证明 $f^{-1}$ 一般不正合.

\textit{提示:} 考虑 $f=i_x$ 是剩余域嵌入的情形, 并使用 (c).
\end{enumerate}
\end{exercise}
习题~\bookref{7.1.4}(b) 表明, 对代数叠的态射 $f:\mathcal X\to\mathcal Y$, 有 $f^{-1}\mathcal O_{\mathcal Y}=\mathcal O_{\mathcal X}$. 因此 $\mathcal O_{\mathcal Y}$-模 $G$ 的拉回是 $\mathcal O_{\mathcal X}$-模 $f^*G=f^{-1}G\otimes_{f^{-1}\mathcal O_{\mathcal Y}}\mathcal O_{\mathcal X}=f^{-1}G$.
\begin{caution}[缺乏函子性]\label{src:7.1.5}
$f^{-1}$ 不正合与光滑平展景 $\mathcal X_{\text{lis-ét}}$ 中缺乏纤维积有关: 若 $U\to V$ 和 $U'\to V$ 是光滑 $\mathcal X$-概形之间的任意态射, 则 $U\times_V U'$ 未必在 $\mathcal X$ 上光滑. 由于没有纤维积, \bookeqref{7.1.2} 中余极限的指标范畴不是滤过的. 指标范畴 $I$ 是滤过的, 当且仅当以 $I$ 为指标的余极限与有限极限交换 \cite[Tag 002W]{SP}, 因而 $f^{-1}$ 不与有限极限交换, 等价地说, 它不左正合. 另见 \cite[Warning 5.3.12]{Beh03}, \cite[例 3.4]{Ols07} 和 \cite[Tag 07BF]{SP}.

遗憾的是, 这意味着光滑平展景不具有函子性: 态射 $f:\mathcal X\to\mathcal Y$ 未必诱导景的态射 $\mathcal X_{\text{lis-ét}}\to\mathcal Y_{\text{lis-ét}}$, 因为景态射的定义 \cite[Tag 00X0]{SP} 要求给出一个连续函子 $\mathcal Y_{\text{lis-ét}}\to\mathcal X_{\text{lis-ét}}$ (方向相反!), 即 $\mathcal Y_{\text{lis-ét}}$ 中覆盖的逆像是 $\mathcal X_{\text{lis-ét}}$ 中的覆盖 (这里确实成立), 并要求层范畴上的逆像函子 $f^{-1}:\mathsf{Sh}(\mathcal Y_{\text{lis-ét}})\to\mathsf{Sh}(\mathcal X_{\text{lis-ét}})$ 正合 (这里不成立). 由于缺乏函子性, 光滑平展景中的构造有时需要格外小心, 尤其是导出范畴和余切复形; 见 \cite{Ols07}. 尽管如此, 光滑平展景中的拟凝聚层 (本书只需用到这些层) 可以很快引入, 并具有与概形情形相同的熟悉性质.
\end{caution}
\begin{remark}[大 fppf 拓扑]\label{src:7.1.6}
可以在大 fppf 拓扑中定义代数叠上的层. 这给出不同的层范畴, 但给出等价的拟凝聚层范畴, 理由与 Deligne--Mumford 叠情形相同 (见习题~\bookref{5.1.25}).

大 fppf 拓扑具有函子性, 但有其他技术问题. Stacks Project 采用这一方法, 因此我们从概形情形开始, 简要概述其主要特点. 概形态射 $f:X\to Y$ 诱导大 fppf 景之间的连续函子 $(\mathsf{Sch}/Y)_{\mathrm{fppf}}\to(\mathsf{Sch}/X)_{\mathrm{fppf}}$, 定义为 $(V\to Y)\mapsto(V\times_Y X\to X)$. 推前与以前一样定义为 $f_*F(V\to Y)=F(V\times_Y X\to X)$, 拉回则由公式 $f^{-1}G(U\to X)=F(U\to X\to Y)$ 给出. 拉回公式之所以简单, 是因为我们在大景中工作: 恒等映射 $U\to V=U$ 是定义预层 $(f^{-1}G)^{\mathrm{pre}}$ 的 \bookeqref{7.1.2} 的指标范畴中的始对象, 而且 $(f^{-1}G)^{\mathrm{pre}}=f^{-1}G$ 已经是层. 一般而言, 与纤维积及终对象都交换的连续函子 $\mathcal C\to\mathcal D$ 诱导正合逆像函子 $\mathsf{Sh}(\mathcal C)\to\mathsf{Sh}(\mathcal D)$, 从而定义景的态射 $\mathcal D\to\mathcal C$ \cite[Tag 00X6]{SP}. 将这一事实用于这里, 可知 $f^{-1}$ 正合, 因而 $(\mathsf{Sch}/Y)_{\mathrm{fppf}}\to(\mathsf{Sch}/X)_{\mathrm{fppf}}$ 定义景的态射 $(\mathsf{Sch}/X)_{\mathrm{fppf}}\to(\mathsf{Sch}/Y)_{\mathrm{fppf}}$.\footnote{函子 $v:(\mathsf{Sch}/Y)_{\mathrm{fppf}}\to(\mathsf{Sch}/X)_{\mathrm{fppf}}$ 有连续左伴随 $u:(\mathsf{Sch}/X)_{\mathrm{fppf}}\to(\mathsf{Sch}/Y)_{\text{lis-ét}}$, 定义为 $(U\to X)\mapsto(U\to X\to Y)$, 也可用它定义 $f_*$ 与 $f^{-1}$: 在预层层面, $f_*=v^p={}_pu$, 且 $(f^{-1})^{\mathrm{pre}}=v_p=u^p$; ${}_pu,u_p,u^p$ 的一般定义见 \cite[Tags 00VC 与 00XF]{SP}, 它们的关系见 \cite[Tag 09VQ]{SP}. 函子 $u$ 不仅连续, 而且余连续 (景的余连续函子 $u:\mathcal C\to\mathcal D$ 要求对每个对象 $c\in\mathcal C$ 及覆盖 $\{d_i\to u(c)\}$, 存在覆盖 $\{c_i\to c\}$, 使覆盖 $\{u(c_i)\to u(c)\}$ 加细 $\{d_i\to u(c)\}$). 这也自动蕴涵 $f^{-1}$ 正合 \cite[Tag 06NW]{SP}.}

对代数叠态射 $f:\mathcal X\to\mathcal Y$, 有类似情形. 大 fppf 景 $\mathcal X_{\mathrm{fppf}}$ 是 $\mathcal X$ 上所有概形组成的范畴 $\mathcal X$, 覆盖取 fppf 覆盖. 此时拉回 $f^{-1}G(U\to\mathcal X)=G(U\to\mathcal X\to\mathcal Y)$ 如上定义 (不必层化), 推前可定义为 $f_*F(V\to\mathcal Y)=\lim_{U\to V}F(U\to\mathcal X)$, 其中极限的指标范畴是 $\mathcal X\to\mathcal Y$ 上概形态射 $U\to V$ 的范畴; 若 $f$ 可由概形表示, 则 $V\times_{\mathcal Y}\mathcal X\to V$ 是终对象, 且 $f_*F(V\to\mathcal Y)=F(V\times_{\mathcal Y}\mathcal X\to\mathcal X)$. 或者, 可以通过光滑下降定义 $F$ 在代数叠的任意态射 $\mathcal U\to\mathcal Y$ 上的截面, 再用通常公式 $f_*F(V\to\mathcal Y)=F(V\times_{\mathcal Y}\mathcal X\to\mathcal X)$ 定义推前. 如上述脚注那样, 用景的连续且余连续函子 $f:\mathcal X_{\mathrm{fppf}}\to\mathcal Y_{\mathrm{fppf}}$ 解释 $f_*$ 和 $f^{-1}$, 可知 $f^{-1}$ 正合. 这意味着 $(f_*,f^{-1})$ 定义景的态射 $\mathcal X_{\mathrm{fppf}}\to\mathcal Y_{\mathrm{fppf}}$, 因而代数叠的大 fppf 景具有函子性 (与光滑平展景不同). 但稍后会看到, $\mathcal O$-模上的推前函子不保持拟凝聚性; 见注记~\bookref{7.1.9}.
\end{remark}
\subsection{拟凝聚层}\label{src:7.1.2-subsection}
\begin{definition}\label{src:7.1.7}
设 $\mathcal X$ 是代数叠. 若 $\mathcal O_{\mathcal X}$-模 $F$ 满足下列条件, 则称它为\emph{拟凝聚的}.
\begin{enumerate}[(1)]
\item 对每个光滑 $\mathcal X$-概形 $U$, 在 $U$ 的小 Zariski 景上的限制 $F|_{U_{\mathrm{Zar}}}$ 是拟凝聚 $\mathcal O_{U_{\mathrm{Zar}}}$-模.
\item 对光滑 $\mathcal X$-概形之间的每个态射 $f:U\to V$, 自然态射 $f^*(F|_{V_{\mathrm{Zar}}})\to F|_{U_{\mathrm{Zar}}}$ 是同构.
\end{enumerate}
若对每个光滑 $\mathcal X$-概形 $U$, $F|_{U_{\mathrm{Zar}}}$ 是向量丛 (分别地, 秩 $r$ 的向量丛, 线丛), 则称 $\mathcal X$ 上的拟凝聚层 $F$ 为\emph{向量丛} (分别地, \emph{秩 $r$ 的向量丛}, \emph{线丛}). 若 $\mathcal X$ 还局部 Noether, 且对每个光滑 $\mathcal X$-概形 $U$, $F|_{U_{\mathrm{Zar}}}$ 凝聚, 则称 $F$ 为\emph{凝聚的}.
\end{definition}
等价地, $\mathcal O_{\mathcal X}$-模 $F$ 拟凝聚, 当且仅当它作为赋环景 $(\mathcal X_{\text{lis-ét}},\mathcal O_{\mathcal X})$ 中的模拟凝聚, 如习题~\bookref{5.1.22}(iii) 所述. 若 $\mathcal X$ 是概形 (分别地, Deligne--Mumford 叠), 则习题~\bookref{5.1.25} 表明上述定义与定义~\bookref{5.1.17} 相容. 不难看到, 代数叠的光滑平展景中的拟凝聚性也等价于大 fppf 拓扑中的拟凝聚性. 无论采用哪种方式, 都得到良定义的拟凝聚层范畴 $\QCoh(\mathcal X)$ 及 (在 Noether 情形下的) 凝聚层范畴 $\mathsf{Coh}(\mathcal X)$. 拟凝聚层的张量积 $F\otimes G:=F\otimes_{\mathcal O_{\mathcal X}}G$ 和 Hom 层 $\mathcal H\!om_{\mathcal O_{\mathcal X}}(F,G)$ 按通常方式定义, 它们是拟凝聚层, 并具有熟悉的性质.
给定代数叠态射 $f:\mathcal X\to\mathcal Y$, 拉回 $f^*=f^{-1}$ 保持拟凝聚性. 若 $f$ 拟紧且拟分离, 则推前 $f_*$ 也保持拟凝聚性, 并是 $f^*$ 的右伴随:
\[
\begin{tikzcd}[ampersand replacement=\&,column sep=large]
\QCoh(\mathcal X)\arrow[r,shift left,"{f_*}"]\&\QCoh(\mathcal Y)\arrow[l,shift left,"{f^*}"] .
\end{tikzcd}
\]
\begin{exercise}[细节]\label{src:7.1.8}
验证 $f_*$ 与 $f^*$ 保持拟凝聚性, 且互为伴随函子.
\end{exercise}
\begin{remark}[大 fppf 景中拟凝聚层的推前]\label{src:7.1.9}
给定代数叠态射 $f:\mathcal X\to\mathcal Y$, 注记~\bookref{7.1.6} 中定义的大 fppf 景中的层推前 $f_*$ 不保持拟凝聚性. 即便对概形态射 $f:X\to Y$, 这一性质也不成立. 若 $F\in\QCoh(X)$, 且
\[
\begin{tikzcd}[ampersand replacement=\&]
U\arrow[r,"{f'}"]\arrow[d,"{g'}"']\&V\arrow[d,"{g}"]\\
X\arrow[r,"{f}"']\arrow[ur,phantom,"{\square}"]\&Y
\end{tikzcd}
\]
是概形的 Cartesian 图, 则由大景中的定义可得 $g^*f_*F=g^{-1}f_*F=f'_*g'^{-1}F=f'_*g'^*F$, 但这一公式一般不成立 (例如 $g$ 平坦时才成立). 换言之, $f_*F$ 只有在平坦态射 $V\to Y$ 上, 才具有所需的拟凝聚推前 $f_{\mathrm{QCoh},*}$ 的正确截面. 稍加小心 (见 \cite[Tag 070A]{SP}), 可将拟凝聚推前定义为复合
\[
\begin{aligned}
&f_{\mathrm{QCoh},*}:\QCoh(\mathcal X)\longrightarrow\operatorname{LQCoh}^{\mathrm{fbc}}(\mathcal O_{\mathcal X})\\
&\hspace{5em}\xrightarrow{f_*}\operatorname{LQCoh}^{\mathrm{fbc}}(\mathcal O_{\mathcal Y})\xrightarrow{Q}\QCoh(\mathcal Y),
\end{aligned}
\]
其中 $\operatorname{LQCoh}^{\mathrm{fbc}}(\mathcal O_{\mathcal X})$ 表示具有平坦基变换性质的局部拟凝聚 $\mathcal O_{\mathcal X}$-模的满子范畴. 局部拟凝聚意指: 对每个来自概形的平坦态射 $U\to\mathcal X$, 限制 $F|_{U_{\text{ét}}}$ 拟凝聚; 平坦基变换性质意指: 对 $\mathcal X$ 上每个平坦态射 $f:U\to V$, 诱导映射 $f^*(F|_{V_{\text{ét}}})\to F|_{U_{\text{ét}}}$ 是同构. $Q$ 是拟凝聚化函子. 函子 $f_{\mathrm{QCoh},*}$ 是 $f^{-1}=f^*$ 的右伴随. 关于用 fppf 拓扑建立导出范畴, 另见 \cite[Tag 08MW]{SP}.
\end{remark}
\begin{exercise}\label{src:7.1.10}
设 $X$ 是概形, $\mathbb G_m$ 是 $\Spec\mathbb Z$ 上的乘法群. 证明 $X\times B\mathbb G_m$ 上的拟凝聚层与 $X$ 上的 $\mathbb Z$-分次拟凝聚层是相同的数据.

\textit{提示:} 推广命题~\bookref{B.1.16} 的论证.
\end{exercise}
\begin{exercise}\label{src:7.1.11}
设 $G$ 是域 $\Bbbk$ 上的仿射代数群. 回顾, $G$-表示是 $\Bbbk$-向量空间, 配备满足两个自然相容条件的余作用 $\sigma:V\to\Gamma(G,\mathcal O_G)\otimes_{\Bbbk}V$ (见 \S\hyperref[src:B.1.13]{B.1.13}).
\begin{enumerate}[(a)]
\item 证明 $\QCoh(BG)$ 等价于 $G$-表示范畴 $\operatorname{Rep}(G)$. 更一般地, 若 $\Spec A$ 是带 $G$-作用的仿射 $\Bbbk$-概形, 证明 $[\Spec A/G]$ 上的拟凝聚层是 $A$-模 $M$ 配备 $\Bbbk$ 上的余作用 $\sigma:M\to\Gamma(G,\mathcal O_G)\otimes_{\Bbbk}M$ 的数据 (即使 $M$ 成为 $G$-表示的 $\Bbbk$-向量空间映射), 满足乘法 $A\otimes_{\Bbbk}M\to M$ 是 $G$-表示的映射. 这推广了 $G$ 有限时的例~\bookref{5.1.21}.
\item 考虑图
\[
\begin{tikzcd}[ampersand replacement=\&]
\Spec A\arrow[r,"{p}"]\&{[\Spec A/G]}\arrow[r,"{\pi}"]\arrow[d,"{q}"]\&\Spec A^G\\
\&BG.\&
\end{tikzcd}
\]
描述拟凝聚层上的函子 $p_*,p^*,\pi_*,\pi^*,q_*,q^*$, 从而推广习题~\bookref{5.1.27}.
\item 若 $U$ 是带 $G$-作用的 $\Bbbk$-概形, 则\emph{带 $G$-作用的线丛} 是 $U$ 上的线丛 $L$ 配备同构 $\alpha:\sigma^*L\xrightarrow{\sim}p_2^*L$, 满足余圈条件 $p_{23}^*\alpha\circ(\mathrm{id}_G\times\sigma)^*\alpha=(\mu\times\mathrm{id}_U)^*\alpha$; 见 \bookref{B.1.31}. 证明带 $G$-作用的线丛与商叠 $[U/G]$ 上的线丛相同.
\end{enumerate}
\end{exercise}
\begin{example}\label{src:7.1.12}
域 $\Bbbk$ 上两个不同构的仿射代数群 $G,H$ 可能具有同构的分类叠 $BG\cong BH$, 从而具有等价的表示范畴. 例如, 若 $O(q),O(q')$ 是关于 $n$ 维 $\Bbbk$-向量空间 $V$ 上非退化二次型 $q,q'$ 的正交群, 则 $BO(q)\cong BO(q')$ (见习题~\bookref{4.1.13}), 因而 $O(q)$ 与 $O(q')$ 具有等价的表示范畴.
\end{example}
我们已经见过 Deligne--Mumford 叠上拟凝聚层的例子, 其中 $\overline{\mathcal M}_g$ 上的多重典范线丛 $\lambda=\det\pi_*(\omega_{\mathcal U_g/\overline{\mathcal M}_g}^{\otimes k})$ 在 $\overline M_g$ 的射影性中起了重要作用 (见 \S\bookref{6.9}). 行列式线丛在向量丛模叠的研究中起类似作用; 见 \S\bookref{9.4.3-subsection}.
\begin{example}\label{src:7.1.13}
若 $\mathcal X$ 是概形 $S$ 上有限呈示的代数叠, 则 $\mathcal X_{\text{lis-ét}}$ 上的相对微分层 $\Omega_{\mathcal X/S}$, 在光滑 $\mathcal X$-概形 $U$ 上定义为 $\Omega_{\mathcal X/S}(U)=\Omega_{U/S}$, 不是拟凝聚的. 因为对光滑 $\mathcal X$-概形之间的非平展态射 $f:U\to V$, $f^*\Omega_{V/S}\to\Omega_{U/S}$ 未必是同构. 这与 Deligne--Mumford 情形不同, 后者中 $\mathcal X_{\text{ét}}$ 上的层 $\Omega_{\mathcal X/S}$ 拟凝聚 (例~\bookref{5.1.19}). 当 $\mathcal X$ 是 Deligne--Mumford 叠时, 可将 $\Omega_{\mathcal X/S}$ 延伸为 $\mathcal X_{\text{lis-ét}}$ 上的拟凝聚层: 对来自概形的光滑态射 $f:U\to\mathcal X$, 将 $\Omega_{\mathcal X/S}(U)$ 定义为 $U_{\text{lis-ét}}$ 上层 $f^*\Omega_{\mathcal X/S}$ 的整体截面.
\end{example}
习题~\bookref{5.1.28} 和 \bookref{5.1.43} 可推广到代数叠.
\begin{proposition}[平坦基变换]\label{src:7.1.14}
考虑代数叠的 Cartesian 图
\[
\begin{tikzcd}[ampersand replacement=\&]
X'\arrow[r,"{g'}"]\arrow[d,"{f'}"']\&X\arrow[d,"{f}"]\\
Y'\arrow[r,"{g}"']\arrow[ur,phantom,"{\square}"]\&Y
\end{tikzcd}
\]
并设 $F$ 是 $X$ 上的拟凝聚层. 若 $g:Y'\to Y$ 平坦, $f:X\to Y$ 拟紧且拟分离, 则自然伴随映射 $g^*f_*F\to f'_*g'^*F$ 是同构.
\end{proposition}
\begin{proposition}\label{src:7.1.15}
设 $\mathcal X$ 是 Noether 代数叠. $\mathcal X$ 上每个拟凝聚层都是其凝聚子层的有向余极限. 若 $\mathcal U\subseteq\mathcal X$ 是开子叠, 则 $\mathcal U$ 上每个凝聚层都可延伸为 $\mathcal X$ 上的凝聚层.
\end{proposition}
\begin{exercise}\label{src:7.1.16}
设 $\mathcal X\to\mathcal Y$ 是具有仿射对角的 Noether 代数叠之间的光滑仿射态射.
\begin{enumerate}[(1)]
\item 证明 $\mathcal X$ 上存在向量丛 $\Omega_{\mathcal X/\mathcal Y}$, 满足: 若 $V\to\mathcal Y$ 是来自概形的态射, 则 $\Omega_{\mathcal X/\mathcal Y}$ 在 $X_V:=\mathcal X\times_{\mathcal Y}V$ 上的拉回是 $\Omega_{X_V/V}$.
\item 给定交换图
\[
\begin{tikzcd}[ampersand replacement=\&]
\Spec A_0\arrow[r,"{f_0}"]\arrow[d,hook]\&\mathcal X\arrow[d]\\
\Spec A\arrow[r]\arrow[ur,dashed]\&\mathcal Y,
\end{tikzcd}
\]
其中 $A\twoheadrightarrow A_0$ 是 Noether 环的满射, 核 $J$ 平方为零. 证明提升的集合是 $\operatorname{Hom}_{A_0}(f_0^*\Omega_{\mathcal X/\mathcal Y},J)$ 下的挠子, 特别地, 它非空.
\item 能否减弱假设?
\end{enumerate}
\end{exercise}
\subsection{拟凝聚构造}\label{src:7.1.3-subsection}
将 \S\bookref{5.1.5-subsection} 中 Deligne--Mumford 叠上的构造推广到代数叠 $\mathcal X$: 拟凝聚 $\mathcal O_{\mathcal X}$-代数是具有相容环对象结构的拟凝聚 $\mathcal O_{\mathcal X}$-模 $\mathcal A$. 相对谱 $\mathcal S\!pec_{\mathcal X}\mathcal A$ 定义为对 $(f,\alpha)$ 组成的叠, 其中 $f:S\to\mathcal X$ 是来自概形的态射, $\alpha:f^*\mathcal A\to\mathcal O_S$ 是 $\mathcal O_S$-代数映射; 它是 $\mathcal X$ 上的仿射代数叠. 在 Noether 代数叠上, 每个拟凝聚 $\mathcal O_{\mathcal X}$-代数都是有限型子代数的有向余极限.
$\mathcal X$ 的约化是 $\mathcal X_{\mathrm{red}}:=\mathcal S\!pec_{\mathcal X}\mathcal O_{\mathcal X}^{\mathrm{red}}$, 其中 $\mathcal O_{\mathcal X}^{\mathrm{red}}$ 是 $\mathcal O_{\mathcal X}$-代数层, 在光滑 $\mathcal X$-概形 $U$ 上定义为 $\mathcal O_{\mathcal X}^{\mathrm{red}}(U)=\Gamma(U,\mathcal O_U)_{\mathrm{red}}$. 若 $\mathcal X$ 整, 则 $\mathcal X$ 的正规化定义为 $\widetilde{\mathcal X}:=\mathcal S\!pec_{\mathcal X}\mathcal A$, 其中 $\mathcal A$ 是 $\mathcal O_{\mathcal X}$-代数, 在光滑 $\mathcal X$-概形 $U$ 上的截面环是 $\Gamma(U,\mathcal O_U)$ 的正规化; 由于正规化与光滑基变换交换 (命题~\bookref{A.7.4}), 这一定义良好.
对代数叠的拟紧且拟分离态射 $f:\mathcal X\to\mathcal Y$, 有分解 $f:\mathcal X\to\mathcal S\!pec f_*\mathcal O_{\mathcal X}\to\mathcal Y$. 态射 $f$ 仿射, 当且仅当 $\mathcal X\to\mathcal S\!pec f_*\mathcal O_{\mathcal X}$ 是同构; 它拟仿射, 当且仅当 $\mathcal X\to\mathcal S\!pec f_*\mathcal O_{\mathcal X}$ 是开浸入. Deligne--Mumford 叠的 Zariski 主定理 (\bookref{5.5.9}) 的证明可推广到代数叠.
\begin{theorem}[Zariski 主定理]\label{src:7.1.17}
Noether 代数叠之间可表示, 拟有限且分离的态射 $f:\mathcal X\to\mathcal Y$ 可分解为稠密开浸入 $\mathcal X\hookrightarrow\widetilde{\mathcal Y}$ 与有限态射 $\widetilde{\mathcal Y}\to\mathcal X$ 的复合.
\end{theorem}
\subsection{Picard 群}\label{src:7.1.4-subsection}
若 $\mathcal X$ 是代数叠, 以 $\Pic(\mathcal X)$ 记 $\mathcal X$ 上线丛的同构类集合. 它在张量积下成为 Abel 群.
\begin{example}\label{src:7.1.18}
若 $G$ 是域 $\Bbbk$ 上的仿射代数群, 则 $\Pic(BG)$ 可与 $G$ 的特征群 $\mathbb X^*(G)$ 等同. 例如, $\Pic(B\mathbb G_m)=\mathbb Z$, $\Pic(B\mu_n)$, $\Pic(B\mathrm{GL}_n)=\mathbb Z$, 且 $\Pic(\mathrm{PGL}_n)=0$. 更一般地, 若 $U$ 是带 $G$-作用的概形, 则 $\Pic([U/G])$ 可与 $U$ 上带 $G$-作用的线丛组成的群 $\Pic^G(U)$ 等同 (习题~\bookref{7.1.11}(c)).
\end{example}
\begin{exercise}\label{src:7.1.19}
设 $\mathcal X$ 是域 $\Bbbk$ 上光滑且不可约的代数叠. 若 $\mathcal D\subseteq\mathcal X$ 是余维 $1$ 的约化子叠, 补为 $\mathcal U$, 证明存在自然定义的线丛 $\mathcal O(\mathcal D)$ (推广概形的通常构造), 使 $\mathcal O(\mathcal D)|_{\mathcal U}\cong\mathcal O_{\mathcal U}$.
\end{exercise}
反之不成立: 线丛未必能表示为 Weil 除子或 Cartier 除子, 例如 $BG$ 上对应于有限群 $G$ 的非平凡特征 $\chi$ 的线丛 $L_\xi$, 或 $\overline{\mathcal M}_{1,1}$ 上的 Hodge 线丛 $\lambda$ (见习题~\bookref{7.1.21}(b)). 为建立良好的 Weil 除子理论, 或更一般的 Chow 群理论, 需要更灵活的闭链概念; 见 \S\bookref{7.1.6-subsection}.
\begin{exercise}\label{src:7.1.20}
设 $\mathbb G_m$ 以权 $d_1,\ldots,d_n$ 作用于域 $\Bbbk$ 上的 $\mathbb A^n$. 设 $\mathcal O(1)$ 是 $[\mathbb A^n/\mathbb G_m]$ 上对应于投影 $[\mathbb A^n/\mathbb G_m]\to B\mathbb G_m$ 的线丛.
\begin{enumerate}[(a)]
\item 证明 $\Pic([\mathbb A^n/\mathbb G_m])\cong\mathbb Z$, 由 $\mathcal O(1)$ 生成.
\item 证明限制 $\Pic([\mathbb A^n/\mathbb G_m])\to\Pic(\mathcal P(d_1,\ldots,d_n))$ 是同构, 其中 $\mathcal P(d_1,\ldots,d_n)$ 是加权射影叠 (见例~\bookref{4.9.14}).
\item 若 $f\in\Gamma(\mathbb A^n,\mathcal O_{\mathbb A^n})$ 是次数 $d$ 的齐次多项式, 且 $V(f)\subseteq\mathbb A^n$ 约化, 证明在 $[\mathbb A^n/\mathbb G_m]$ 上有 $\mathcal O(V(f))\cong\mathcal O(d)$.
\end{enumerate}
\end{exercise}
下面给出 Mumford 的计算 $\Pic(\mathcal M_{1,1})=\mathbb Z/12$ 的初等方法, 这一计算首先作为习题~\bookref{5.1.50} 中的困难习题提出.
\begin{exercise}\label{src:7.1.21}
设 $k$ 是域, $\operatorname{char}(k)\neq2,3$.
\begin{enumerate}[(a)]
\item 证明 $\Pic(\overline{\mathcal M}_{1,1})=\mathbb Z\langle\lambda\rangle$.

\textit{提示:} 使用习题~\bookref{4.1.21}(b) 中的描述 $\overline{\mathcal M}_{1,1}=[(\mathbb A^2\setminus0)/\mathbb G_m]$, 其中 $\mathbb G_m$ 以权 $4,6$ 作用. 证明限制 $\Pic([\mathbb A^2/\mathbb G_m])\to\Pic(\overline{\mathcal M}_{1,1})$ 是等价.
\item 证明在 $\overline{\mathcal M}_{1,1}$ 上有 $12\lambda=\delta$, 且 $12\lambda$ 是 $\lambda$ 能下降到 $\overline M_{1,1}$ 的最小倍数. 证明不存在有效 Cartier 除子 $\mathcal D\subset\overline{\mathcal M}_{1,1}$ 使 $\lambda=\mathcal O_{\overline{\mathcal M}_{1,1}}(\mathcal D)$.
\item 证明 $\Pic(\mathcal M_{1,1})=\mathbb Z/12$.

证明限制 $\Pic(\overline{\mathcal M}_{1,1})\to\Pic(\mathcal M_{1,1})$ 满射, 且 $\mathcal O(\Delta)=\mathcal O(12)$ 的像平凡. 考虑 $\mathcal O(4)$ 和 $\mathcal O(6)$ 在自同构群为 $\mathbb Z/4$ 和 $\mathbb Z/6$ 的唯一椭圆曲线的剩余 gerbe 上的限制, 证明它们的像非平凡. 另见 \cite{Mum65} (任意基底上的推广见 \cite{FO10}).
\end{enumerate}
\end{exercise}
\begin{exercise}\label{src:7.1.22}
设 $X$ 是概形, $L$ 是 $X$ 上的线丛, $r>0$. 考虑 $r$ 次根 gerbe $X(\sqrt[r]{L})$ (见例~\bookref{4.9.21} 和习题~\bookref{7.4.26}). 证明
\[
\Pic(X(\sqrt[r]{L}))\cong\Pic(X)[\sqrt[r]{L}],
\]
其中 $\Pic(X)[\sqrt[r]{L}]$ 是 $\Pic(X)[z]/(rz-[L])$ 的简写.
\end{exercise}
\subsection{层上同调}\label{src:7.1.5-subsection}
代数叠的层上同调可用与 \S\bookref{5.1.6-subsection} 中 Deligne--Mumford 叠情形基本相同的方法建立.
\begin{lemma}\label{src:7.1.23}
若 $\mathcal X$ 是代数叠, 则范畴 $\mathsf{Ab}(\mathcal X_{\text{lis-ét}})$ 和 $\QCoh(\mathcal X)$ 有足够多内射对象.
\end{lemma}
\begin{proof}
任意景上的 Abel 层范畴有足够多内射对象 \cite[Tag 01DP]{SP}, 而 $\QCoh(\mathcal X)$ 是 Grothendieck Abel 范畴, 因而有足够多内射对象 \cite[Tag 079H]{SP}.
\end{proof}
\begin{definition}[上同调]\label{src:7.1.24}
设 $\mathcal X$ 是代数叠, $F$ 是 $\mathcal X_{\text{lis-ét}}$ 上的 Abel 群层. 上同调群 $\mathrm H^i(\mathcal X_{\text{lis-ét}},F)$ 定义为整体截面函子 $\Gamma:\mathsf{Ab}(\mathcal X_{\text{lis-ét}})\to\mathsf{Ab}$ 的第 $i$ 个右导出函子.
\end{definition}
Čech 上同调的定义与 Deligne--Mumford 情形 (定义~\bookref{5.1.37}) 相同.
\begin{definition}[Čech 上同调]\label{src:7.1.25}
给定代数叠的光滑覆盖 $\mathcal U=\{U_i\to\mathcal X\}_{i\in I}$ 及 $\mathcal X_{\text{lis-ét}}$ 上的 Abel 层 $F$, $F$ 关于 $\mathcal U$ 的 Čech 复形为 $\check{\mathcal C}^{\bullet}(\mathcal U,F)$, 其中
\[
\check{\mathcal C}^n(\mathcal U,F)=\prod_{(i_0,\ldots,i_n)\in I^{n+1}}F(U_{i_0}\times_{\mathcal X}\cdots\times_{\mathcal X}U_{i_n}),
\]
微分为
\[
\begin{aligned}
&d^n:\check{\mathcal C}^n(\mathcal U,F)\longrightarrow\check{\mathcal C}^{n+1}(\mathcal U,F),\\
&(s_{i_0,\ldots,i_n})\longmapsto\biggl(\sum_{k=0}^{n+1}(-1)^kp_{\widehat k}^*s_{i_0,\ldots,\widehat{i_k},\ldots,i_n}\biggr)_{(i_0,\ldots,i_{n+1})},
\end{aligned}
\]
其中 $p_{\widehat k}:U_{i_0}\times_{\mathcal X}\cdots\times_{\mathcal X}U_{i_n}\to U_{i_0}\times_{\mathcal X}\cdots\times_{\mathcal X}\widehat{U_{i_k}}\times_{\mathcal X}\cdots\times_{\mathcal X}U_{i_n}$ 是忘却第 $k$ 个分量的映射 (指标从 $0$ 开始). $F$ 关于 $\mathcal U$ 的 Čech 上同调为
\[
\check{\mathrm H}^i(\mathcal U,F):=\mathrm H^i(\check{\mathcal C}^{\bullet}(\mathcal U,F)).
\]
\end{definition}
\begin{theorem}\label{src:7.1.26}
对仿射概形 $X$ 上的拟凝聚 $\mathcal O_{X_{\text{lis-ét}}}$-模 $F$, 对所有 $i>0$ 都有 $\mathrm H^i(X_{\text{lis-ét}},F)=0$.
\end{theorem}
\begin{proof}
小平展拓扑中 $\mathrm H^i(X_{\text{ét}},F)=0$ 的论证 (定理~\bookref{5.1.40}) 可推广.
\end{proof}
若对所有 $p>0,n\geq0$ 及 $i_0,\ldots,i_n\in I$ 都有 $\mathrm H^p(U_{i_0}\times_U\cdots\times_U U_{i_n},F)=0$, 则由引理~\bookref{5.1.39} 中相同的形式论证可得 $\mathrm H^p(\mathcal X_{\text{ét}},F)=\check{\mathrm H}^p(\mathcal U,F)$. 特别地, 得到:
\begin{lemma}\label{src:7.1.27}
设 $\mathcal X$ 是具有仿射对角的代数叠, $F$ 是拟凝聚层. 若 $\mathcal U=\{U_i\to\mathcal X\}$ 是平展覆盖, 每个 $U_i$ 都仿射, 则 $\mathrm H^i(\mathcal X_{\text{lis-ét}},F)=\check{\mathrm H}^i(\mathcal U,F)$.
\end{lemma}
\begin{proposition}\label{src:7.1.28}
若 $\mathcal X$ 是具有仿射对角的 Deligne--Mumford 叠 (例如概形), $F$ 是拟凝聚层, 则对所有 $i$, $\mathrm H^i(\mathcal X_{\text{ét}},F)=\mathrm H^i(\mathcal X_{\text{lis-ét}},F)$.
\end{proposition}
\begin{proof}
命题~\bookref{5.1.41} 的论证可推广.
\end{proof}
类似论证表明, 拟凝聚层 $F$ 的上同调也可在大 fppf 拓扑中计算. 因此通常省略下标, 简记为 $\mathrm H^i(\mathcal X,F)$.
\begin{proposition}\label{src:7.1.29}
若 $\mathcal X$ 是具有仿射对角的代数叠, $F$ 是拟凝聚层, 则 $F$ 作为 Abel 层的上同调 $\mathrm H^i(\mathcal X_{\text{lis-ét}},F)$ 与 $\Gamma:\QCoh(\mathcal X)\to\mathsf{Ab}$ 的第 $i$ 个右导出函子一致.
\end{proposition}
\begin{proof}
命题~\bookref{5.1.42} 的论证可推广.
\end{proof}
若 $\mathcal X$ 没有仿射对角, 则拟凝聚层的层上同调 $\mathrm H^i(\mathcal X,F)$ 可能与 $\Gamma(\mathcal X,-):\QCoh(\mathcal X)\to\mathsf{Ab}$ 的第 $i$ 个右导出函子不同.
\begin{exercise}\label{src:7.1.30}
设 $E$ 是域上的椭圆曲线, $BE$ 是其分类叠. 证明对每个 $i$, $\mathrm H^{2i}(BE,\mathcal O_{BE})\neq0$, 但若在 $\QCoh(BE)$ 中计算, 则上同调为零.
\end{exercise}
\begin{proposition}\label{src:7.1.31}
若 $\mathcal X$ 是代数叠, $F_i$ 是 $\mathcal X_{\text{lis-ét}}$ 中 Abel 层的有向系统, 则 $\operatorname{colim}_i\mathrm H^i(\mathcal X_{\text{lis-ét}},F_i)\to\mathrm H^i(\mathcal X_{\text{lis-ét}},\operatorname{colim}_iF_i)$ 是同构.
\end{proposition}
\begin{example}[群上同调]\label{src:7.1.32}
域 $\Bbbk$ 上仿射代数群 $G$ 的群上同调可以解释为 $BG$ 上的层上同调, 正如有限群情形 (例~\bookref{5.1.45}) 一样. 事实上, $G$-表示 $V$ 可视为拟凝聚 $\mathcal O_{BG}$-模, 由命题~\bookref{7.1.29}, 层上同调 $\mathrm H^i(BG,V)$ 与群上同调 $\mathrm H^i(G,V)$ 一致, 后者定义为由 $V\mapsto V^G$ 给出的函子 $\operatorname{Rep}(G)\to\mathsf{Vect}_{\Bbbk}$ 的第 $i$ 个导出函子. 而且, 由引理~\bookref{7.1.27}, 可用 $V$ 关于光滑覆盖 $\mathcal U=\{\Spec\Bbbk\to BG\}$ 的 Čech 复形 $\check C^{\bullet}(\mathcal U,V)$ 计算上同调. 这个复形的项为 $\check C^n(\mathcal U,V)=\Gamma(G,\mathcal O_G)^{\otimes n}\otimes V$, 微分由 \bookeqref{5.1.46} 给出, 对应于群上同调中的标准消解或 bar 消解.

若 $G$ 线性约化, 即 $V\mapsto V^G$ 正合, 则对每个 $G$-表示 $V$ 和 $i>0$, $\mathrm H^i(G,V)=0$. 若 $G$ 作用于 $\Bbbk$ 上的仿射概形 $\Spec A$, 则对每个拟凝聚层 $F$ 和 $i>0$, $\mathrm H^i([\Spec A/G],F)=0$.
\end{example}
\begin{exercise}\label{src:7.1.33}
\begin{enumerate}[(a)]
\item 若 $\Bbbk$ 的特征为 $0$, $\Bbbk$ 是平凡 $\mathbb G_a$-表示, 证明 $\mathrm H^i(B\mathbb G_a,\Bbbk)$ 在 $i=0,1$ 时等于 $\Bbbk$, 其他情形为零.
\item 若 $G$ 是域 $\Bbbk$ 上的仿射代数群, 证明 $\mathrm H^i(BG,V)$ 对所有 $i$ 都有限维, 且当 $i\gg0$ 时为零.

\textit{提示:} 选取嵌入 $G\hookrightarrow\mathrm{GL}_n$, 使用 $\mathrm{GL}_n$ 在特征 $0$ 时线性约化 (定理~\bookref{B.1.46}), 并考虑复合 $BG\to B\mathrm{GL}_n\to\Spec\Bbbk$.
\item 证明 $\mathrm H^1(B\mathbb G_a,\Bbbk)$ 无限维.
\end{enumerate}
\end{exercise}
在正特征下, $\mathrm H^i(G,V)$ 可对任意大的 $i$ 非零; 这甚至出现在 $\mathbb Z/p\mathbb Z$ 等有限群的情形 (习题~\bookref{5.1.47}).
\subsection{Chow 群}\label{src:7.1.6-subsection}
将代数叠上的闭链定义为整闭子叠的形式和 (类似概形情形), 不能得到令人满意的理论. 即便在余维 $1$ 的情形, 闭子叠也不够多; 见习题~\bookref{7.1.19}. 按照 \cite{Tot99,EG98}, 我们改为通过考虑向量丛的闭子概形来定义商叠的 Chow 群.
设 $G$ 是代数闭域 $\Bbbk$ 上维数 $g$ 的光滑仿射代数群, $X$ 是 $\Bbbk$ 上 $n$ 维有限型概形. 对每个 $i$, 选取 $r$ 维 $G$-表示 $V$, 使存在非空开子概形 $U\subseteq\mathbb A(V)$ 满足: (1) $G$ 在 $U$ 上自由作用, (2) 商 $U/G$ 是概形, (3) $\operatorname{codim}(\mathbb A(V)\setminus U)>n-i-g$. 不难证明这样的表示存在. 定义 $X$ 的第 $(i-g)$ 个\newterm{等变 Chow 群}{equivariant Chow group}, 或等价地, $[X/G]$ 的第 $i$ 个 \newterm{Chow 群}{Chow group}, 为
\[
\mathrm{CH}_{i-g}^G(X)=\mathrm{CH}_i([X/G]):=\mathrm{CH}_{i+r}(X\times^G U).
\]
\begin{exercise}\label{src:7.1.34}
\begin{enumerate}[(a)]
\item 证明存在满足性质 (1)--(3) 的表示 $V$.
\item 检查 $\mathrm{CH}_{i-g}^G(X)$ 的定义与选择无关.
\end{enumerate}
\end{exercise}
这一定义乍看可能神秘, 但若要求 Chow 群在高余维开浸入及向量丛下不变, 则它实际上是必然的:
\[
\begin{tikzcd}[ampersand replacement=\&,column sep=small]
X\times^G U\arrow[r,"{\text{开}}"]\&{[(X\times\mathbb A(V))/G]}\arrow[d,"{\text{向量丛}}"]\\
\&{[X/G]}
\end{tikzcd}
\]
\[
\begin{tikzcd}[ampersand replacement=\&,column sep=small]
\mathrm{CH}_{i+r}(X\times^G U)\&\mathrm{CH}_{i+r}([(X\times\mathbb A(V))/G])\arrow[l,"{\sim}"']\\
\&\mathrm{CH}_i([X/G]).\arrow[u,"{\sim}"']
\end{tikzcd}
\]
若 $[X/G]$ 光滑且纯维数为 $d=n-g$, 则定义
\begin{align*}
\mathrm{CH}^i_G(X)&=\mathrm{CH}^i([X/G]):=\mathrm{CH}_{d-i}([X/G]),\\
\mathrm{CH}^*_G(X)&=\mathrm{CH}^*([X/G]):=\bigoplus_i\mathrm{CH}^i([X/G]).
\end{align*}
相交积给出环结构; 称 $\mathrm{CH}^*_G(X)$ 为 $X$ 的\newterm{等变 Chow 环}{equivariant Chow ring}, 称 $\mathrm{CH}^*([X/G])$ 为 $[X/G]$ 的\newterm{Chow 环}{Chow ring}.
\begin{example}[\texorpdfstring{$\mathrm{CH}^*(B\mathbb G_m)$}{CH*(BGm)}]\label{src:7.1.35}
设 $V$ 是 $\Bbbk$ 上权均为 $1$ 的 $r$ 维 $\mathbb G_m$-表示. 则 $\mathbb G_m$ 在 $\mathbb A^n\setminus0$ 上自由作用, 当 $-1\geq i>-r-1$ 时有 $\mathrm{CH}_i(B\mathbb G_m)=\mathrm{CH}_{i+r}(\mathbb P^{r-1})=\mathbb Z$. 因而 $\mathrm{CH}_i(B\mathbb G_m)$ 在 $i\leq-1$ 时为 $\mathbb Z$, 其他情形为 $0$. 于是对 $j\geq0$, $\mathrm{CH}^j(B\mathbb G_m)=\mathbb Z$, 且 $\mathrm{CH}^*(B\mathbb G_m)=\mathbb Z[x]$. 更一般地, 若 $T\cong\mathbb G_m^r$ 是秩 $r$ 的环面, 则 $\mathrm{CH}^*(BT)$ 同构于 $T$ 的特征环 $\mathbb Z[x_1,\ldots,x_r]$.
\end{example}
我们汇总等变 Chow 群的一些重要性质.
\begin{properties}\label{src:7.1.36}
\begin{enumerate}[(1)]
\item (与商呈示无关) 若 $[X/G]\cong[X'/G']$, 则 $\mathrm{CH}_i^G(X)\cong\mathrm{CH}_i^{G'}(X')$; 特别地, $\mathrm{CH}_i([X/G])$ 的定义与商呈示无关.
\item (向量丛不变性) 若 $Y\to X$ 是 $G$-等变且 Zariski 局部为相对维数 $r$ 的仿射纤维化 (例如秩 $r$ 向量丛的总空间), 则 $\mathrm{CH}_*^G(X)\cong\mathrm{CH}_{*+r}^G(Y)$.
\item (切除序列) 若 $\mathcal Z\subseteq\mathcal X=[X/G]$ 是闭子叠, 补为 $\mathcal U$, 则有右正合序列
\[
\mathrm{CH}_*(\mathcal Z)\longrightarrow\mathrm{CH}_*(\mathcal X)\longrightarrow\mathrm{CH}_*(\mathcal U)\longrightarrow0.
\]
\item (与粗模空间的比较) 若 $\mathcal X\cong[X/G]$ 是分离 Deligne--Mumford 叠, 粗模空间为 $X$, 则 $\mathrm{CH}_*(\mathcal X)\otimes\mathbb Q\cong\mathrm{CH}_*(X)\otimes\mathbb Q$.
\item (函子性) 平坦态射诱导 Chow 群上的拉回映射, 固有态射诱导推前映射.
\item (自交) 若 $\mathcal X=[X/G]$ 光滑, $i:\mathcal Z\hookrightarrow\mathcal X$ 是纯余维 $d$ 的光滑子叠, 则与 $\mathcal Z$ 相交给出拉回 $i^*:\mathrm{CH}^*(\mathcal X)\to\mathrm{CH}^*(\mathcal Z)$, 满足对 $\alpha\in\mathrm{CH}^*(\mathcal Z)$,
\[
i^*i_*\alpha=c_d(\mathcal N_{\mathcal Z/\mathcal X})\cap\alpha,
\]
其中 $c_d$ 是法丛的最高 Chern 类.
\item (余维 $1$) 若 $\mathcal X=[X/G]$ 光滑, 则 $\mathrm{CH}^1(\mathcal X)=\Pic(\mathcal X)$.
\end{enumerate}
\end{properties}
\begin{exercise}[困难]\label{src:7.1.37}
证明性质~\bookref{7.1.36}.

\textit{提示:} 归约到概形 Chow 群的通常性质, 参见 \cite{Ful98}.
\end{exercise}
Kresch 将 Chow 群理论推广到 (未必是商叠的) 代数叠; 见 \cite{Kre99}.
\begin{exercise}[困难]\label{src:7.1.38}
\begin{enumerate}[(a)]
\item 若环面 $T$ 在光滑概形 $X$ 上平凡作用, 证明 $\mathrm{CH}_T^*(X)\cong\mathrm{CH}^*(X)\otimes\mathrm{CH}^*(BT)$. 是两个环面的乘积, 且 $T_2$ 平凡作用. 则 $\mathrm{CH}_T^*(X)\cong\mathrm{CH}_{T_1}^*(X)\otimes\mathrm{CH}^*(BT_2)$.
\item 若 $G$ 是具有极大环面 $T$ 的连通约化群, $X$ 是带 $G$-作用的光滑概形, 则 Weyl 群 $W=N_G(T)/T$ 作用于 $\mathrm{CH}_T^*(X)$. 证明 $\mathrm{CH}_G^*(X)_{\mathbb Q}=\mathrm{CH}_T^*(X)_{\mathbb Q}^W$, 且若 $G$ 还特殊 (例~\bookref{B.1.61}), 则 $\mathrm{CH}_G^*(X)=\mathrm{CH}_T^*(X)^W$.
\end{enumerate}
\end{exercise}
\begin{exercise}[射影丛公式]\label{src:7.1.39}
若 $E$ 是光滑纯维数商叠 $\mathcal X=[X/G]$ 上秩 $r+1$ 的向量丛, $G$ 是光滑仿射群概形, 证明
\[
\mathrm{CH}^*(\mathbb P(E))=\mathrm{CH}^*(\mathcal X)[t]/(t^{r+1}+c_1(E)t^r+\cdots+c_{r+1}(E)),
\]
正如概形情形一样 (\cite[定理 3.3]{Ful98}, \cite[定理 9.6]{EH16}).
\end{exercise}
\begin{exercise}[中等]\label{src:7.1.40}
\begin{enumerate}[(a)]
\item 设 $\mathcal P(d_0,\ldots,d_n)$ 是例~\bookref{4.9.14} 的加权射影叠. 证明 $\mathrm{CH}^*(\mathcal P(d_0,\ldots,d_n))\cong\mathbb Z[x]/(d_1\cdots d_nx^{n+1})$.
\item 若 $\operatorname{char}(\Bbbk)\neq2,3$, 证明 $\mathrm{CH}^*(\mathcal M_{1,1})\cong\mathbb Z[x]/(12x)$, 且 $\mathrm{CH}^*(\overline{\mathcal M}_{1,1})\cong\mathbb Z[x]/(24x^2)$. (与习题~\bookref{7.1.21} 的 Picard 群比较.)
\item 设 $\mathbb G_m$ 以权 $d_0,\ldots,d_n$ 作用于 $\mathbb P^n$. 证明 $\mathrm{CH}^*([\mathbb P^n/\mathbb G_m])=\mathbb Z[h,t]/p(h,t)$, 其中 $p(h,t)=\sum_{i=0}^{n+1}h^{n+1-i}e_i(d_0t,\ldots,d_nt)$, $e_i$ 是第 $i$ 个初等对称多项式.
\end{enumerate}
\end{exercise}
\subsection{de Rham 与奇异上同调}\label{src:7.1.7-subsection}
按照 \cite{Beh04}, 我们简要讨论代数叠的 de Rham 与奇异上同调.
\paragraph{解析化.}
若 $\mathcal X$ 是 $\mathbb C$ 上具有仿射对角的光滑代数叠, 则类似有限型 $\mathbb C$-概形的解析化, 它有\newterm{解析化}{analyticification} $\mathcal X^{\mathrm{an}}$, 且 $\mathcal X^{\mathrm{an}}$ 是\newterm{可微叠}{differentiable stack}. 若 $U_0\to\mathcal X$ 是由概形给出的光滑呈示, 因而 $\mathcal X$ 是光滑群胚 $U_1:=U_0\times_{\mathcal X}U_0\rightrightarrows U_0$ 的商, 则 $U_0^{\mathrm{an}}\to\mathcal X^{\mathrm{an}}$ 是光滑呈示, $\mathcal X^{\mathrm{an}}$ 是 Lie 群胚 $U_1^{\mathrm{an}}\rightrightarrows U_0^{\mathrm{an}}$ 的商.
\paragraph{可微叠的 de Rham 上同调.}
给定可微叠 $\mathcal X$ 及其光滑呈示 $U_0\to\mathcal X$, 可定义\newterm{单纯流形}{simplicial manifold}
\begin{equation}\label{src:7.1.41}
\begin{aligned}
&U_\bullet:\begin{tikzcd}[ampersand replacement=\&,column sep=small]
\cdots\arrow[r]\&U_3\arrow[r,shift left=1ex]\arrow[r]\arrow[r,shift right=1ex]\arrow[r,shift right=2ex]\&U_2\arrow[r,shift left=1ex]\arrow[r]\arrow[r,shift right=1ex]\&U_1\arrow[r,shift left=.5ex]\arrow[r,shift right=.5ex]\&U_0,
\end{tikzcd}\\
&\text{其中 }U_p:=\underbrace{U_0\times_{\mathcal X}\cdots\times_{\mathcal X}U_0}_{p+1\text{ 次}},
\end{aligned}
\end{equation}
微分映射 $\partial_0,\ldots,\partial_p:U_p\to U_{p-1}$ 忘却第 $i$ 项, 退化映射 $s_i:U_{p-1}\to U_p$ 在第 $i$ 项插入恒等态射. 这定义了双复形 $\Omega^q(U_p)$, 微分由外微分 $d:\Omega^{q-1}(U_p)\to\Omega^q(U_p)$ 和 $\partial:=\sum_{i=0}^p(-1)^i\partial_i^*:\Omega^q(U_{p-1})\to\Omega^q(U_p)$ 给出. 定义 \newterm{de Rham 复形}{de Rham complex} $C_{\mathrm{dR}}^{\bullet}(\mathcal X)$ 为全复形
\[
C_{\mathrm{dR}}^k(\mathcal X):=\bigoplus_{p+q=k}\Omega^q(U_p),
\]
微分 $\delta:C_{\mathrm{dR}}^k(\mathcal X)\to C_{\mathrm{dR}}^{k+1}(\mathcal X)$ 对 $\omega\in\Omega^p(U_q)$ 定义为 $\delta(\omega)=\partial(\omega)+(-1)^pd(\omega)$. \newterm{de Rham 上同调}{de Rham cohomology} 为
\[
\mathrm H^n_{\mathrm{dR}}(\mathcal X):=\mathrm H^n(C_{\mathrm{dR}}^{\bullet}(\mathcal X)),
\]
它与呈示的选择无关.
\paragraph{拓扑叠的奇异同调与上同调.}
对拓扑叠 $\mathcal X$, 可以重复奇异同调与上同调的构造. 设 $U_0\to\mathcal X$ 是呈示, $U_\bullet$ 是 \bookeqref{7.1.41} 中的单纯拓扑空间. 对每个 $p$, 有奇异链复形 $C_\bullet(U_p)$, 微分为 $d:C_q(U_p)\to C_{q-1}(U_p)$. 利用由映射 $\partial_i:U_p\to U_{p-1}$ 诱导的微分 $\partial=\sum_{i=0}^p(-1)^i\partial_i:C_q(U_p)\to C_q(U_{p-1})$, 得到双复形 $C_q(U_p)$. $\mathcal X$ 的奇异链复形 $C_\bullet(\mathcal X)$ 定义为全复形
\[
C_k(\mathcal X):=\bigoplus_{p+q=k}C_q(U_p),
\]
微分 $\delta:C_k(\mathcal X)\to C_{k-1}(\mathcal X)$ 对 $\gamma\in C_q(U_p)$ 由 $\delta(\gamma)=(-1)^{p+q}\partial(\gamma)+(-1)^qd(\gamma)$ 给出. 对 Abel 群 $A$, 因而可定义 $\mathcal X$ 的 $A$ 系数奇异同调群为
\[
\mathrm H_n(\mathcal X,A):=\mathrm H_n(C_\bullet(\mathcal X)\otimes_{\mathbb Z}A).
\]
对偶化后, 定义奇异上链复形 $C^{\bullet}(\mathcal X)$ 为 $C^n(\mathcal X):=\operatorname{Hom}(C_n(\mathcal X),\mathbb Z)$, $\mathcal X$ 的 $A$ 系数奇异上同调群为
\[
\mathrm H^n(\mathcal X,A):=\mathrm H^n(C^{\bullet}(\mathcal X)\otimes_{\mathbb Z}A).
\]
\paragraph{比较.}
\begin{itemize}
\item 有配对 $\mathrm H_k(\mathcal X,\mathbb Z)\otimes\mathrm H^k(\mathcal X,\mathbb Z)\to\mathbb Z$, 与 $\mathbb Q$ 作张量积后给出等同 $\mathrm H^k(\mathcal X,\mathbb Q)\cong\mathrm H_k(\mathcal X,\mathbb Q)^{\vee}$.
\item 若拓扑群 $G$ 作用于空间 $U$, 则等变上同调定义为 $\mathrm H_G^*(U,A):=\mathrm H^*(EG\times^G U,A)$, 其中 $EG$ 是具有自由 $G$-作用的可缩空间, 且有等同 $\mathrm H^*([U/G],A)=\mathrm H_G^*(U,A)$.
\item 对可微叠 $\mathcal X$, $\mathrm H_{\mathrm{dR}}^*(\mathcal X)$ 可与常值层 $\mathbb R$ 的上同调 $\mathrm H^*(\mathcal X,\mathbb R)$ 等同, 后者在 $\mathcal X$ 上流形的大景中计算, 覆盖取通常的开覆盖.
\item 若 $\mathcal X$ 是拓扑 Deligne--Mumford 叠 (例如 $\mathbb C$ 上分离有限型 Deligne--Mumford 叠所关联的拓扑叠), 粗模空间为 $\mathcal X\to X$, 则 $\mathrm H^*(\mathcal X,\mathbb Q)=\mathrm H^*(X,\mathbb Q)$.
\end{itemize}
